\documentclass[book.tex]{subfiles}
\begin{document}

\chapter{Foundations of Low Data Learning}
\label{chap:low_data}
\noindent\rule{11cm}{0.4pt} \\
\textbf{Prerequisites:}  \autoref{chap:eft},   \\
\textbf{Difficulty Level:} ** \\
\noindent\rule{11cm}{0.4pt}
\vspace{5mm}

"""
The relevant limit for machine learning is not N $\to$ infinity but instead N $\to$ 0. The human visual system is proof that it is possible to learn categories with extremely few samples. This talk will discuss steps towards building such systems. It is structured in three parts. The first part will discuss algorithms to adapt representations of deep networks to new categories with few labeled data. The second part will exploit a formal connection between thermodynamics and machine learning to characterize such adaptation. We will see how this theory leads to algorithms for transfer learning that can guarantee high accuracy on the target task. It stands to reason that adaptation cannot always work, it should be difficult for tasks that are "far away". The third part will characterize when learning multiple tasks is difficult, and demonstrate ways to work around such difficulties.
"""

References \cite{dhillon2019baseline}, \cite{gao2020free}, \cite{gao2021information}, \cite{ramesh2021boosting}

\end{document}